\begin{abstract}
Automated test generation tools for repository-level software maintenance have significantly increased in recent years. These tools excel in synthesizing unit test inputs but are limited by their static context paradigms, which suffer from compliance bias---the tendency to passively generate assertions that conform to existing defective code. This article introduces an open-source, dual-agent framework for proactive repository-level test suite generation through dynamic terminal execution feedback. Our approach pairs an Exploration Agent (DeepSeek-V3 operating via the SWE-agent CLI bash terminal) with a Code Action Fixer Agent (Llama-3.3-70B). By executing candidate test suites in real-time, observing pytest tracebacks, and iteratively refining adversarial assertions across up to 80 turns, our system dynamically rejects compliance assumptions and exposes latent logic bugs. Evaluation results on 1,552 Python tasks from 482 public repositories in the TestExplora benchmark show that interactive terminal feedback achieves an average Pass@1 Fail-to-Pass rate of 35.05\%, more than double the static baseline's 16.60\%. When expanding the sample budget to K=3 independent runs, the cumulative bug discovery rate reaches 71.59\% (paired one-tailed Wilcoxon $p = 0.000000$, Cliff's $\delta = 0.3018$, medium effect size). Furthermore, interactive tracebacks enabled the detection of latent logic faults across 10 distinct software domain categories that passed completely unnoticed by static completion tools.
\end{abstract}

\begin{IEEEkeywords}
Automated Test Generation, Large Language Models, Agentic Exploration, Interactive Terminal Feedback, Compliance Bias, TestExplora Benchmark.
\end{IEEEkeywords}
